\documentclass{situationist-preprint}

\title{Selection as Artistic Agency}
\msid{AES-14-b}
\msfamily{AES-14}
\msdiscipline{Aesthetics}
\mstype{Companion manuscript}
\mskeywords{photography; selection; readymade; collage; film editing; artistic agency}

\begin{document}
\maketitle

\begin{abstract}
AES-14-a argues that photography is art despite its causal transparency, through choices of framing and selection. The weaker of the two is selection: choosing one exposure among many may seem the act of an editor rather than a maker. We argue that selection is a form of artistic agency with a well-established place in the arts, and that photographic selection shares its structure with three practices whose artistic standing is not in doubt: the readymade, collage and film editing. In each, the maker does not fabricate the elements, and in each the aesthetic properties of the work depend on what was chosen, from what range, and in what relation to other elements. The parallel shows that the objection to selection, if sound, would tell equally against these practices.
\end{abstract}

\section{Selection in photography}

By Definition~2.2 of AES-14-a, selection is the choice of one exposure among those the photographer has made. Lemma~3.2 of AES-14-a states that the aesthetic properties of a photograph depend on its framing and selection, and Lemma~3.3 that they are appreciable as the photographer's achievement. The objection we address holds that selection, unlike framing, contributes nothing to the photograph: the selected exposure already had all its properties before it was chosen, and the photographer only recognized them.

The objection assumes that a contribution to a work must be a contribution to the features of its elements. We show that this assumption is false in three established arts.

\section{Three parallels}

\paragraph{The readymade.} When an ordinary manufactured object is presented as a work, its maker fabricates nothing. On the account of AES-01, its standing as art depends on its presentation within an artworld and its critical uptake. What is appreciated is the choice: of this object rather than another, at this moment, in this context of presentation \citep{Danto1964}. The aesthetic and artistic properties ascribed to such works (wit, provocation, austerity) are properties the object has as selected and presented, not properties it had in the hardware store.

\paragraph{Collage.} In collage the elements are found: newspaper, wallpaper, labels, photographs. The work's properties depend on which fragments were chosen and how they were combined. A fragment of newsprint has no compositional role until it is placed; its role is fixed by the choice of it among others and its relation to them.

\paragraph{Film editing.} A film is assembled from shots, many of which are discarded. The choice of take and the order and duration of the shots fix properties of the film, including its rhythm and much of its emotional character, that no single shot has. The early experiments associated with Kuleshov, in which the same shot of a face was read differently depending on what it was cut with, are the standard illustration \citep{Bordwell1985}.

\begin{proposition}\label{p:sel}
In each of the three practices, the selection of elements from a range of alternatives fixes aesthetic properties of the work that the elements, considered apart from the selection, do not have.
\end{proposition}

The same holds of photographic selection, for two reasons. First, the properties of an exposure that count aesthetically, such as the decisiveness of a moment, are comparative: the moment is decisive relative to the ones before and after it that the photographer also captured and rejected \citep{CartierBresson1952}. Second, a selected exposure is often one of a sequence or body of work, and its properties within that body, as with the shots of a film, depend on what else was selected.

\section{Objections}

\begin{objection}
Selection is too easy. Anyone can pick the best of a set of snapshots; if selection makes art, every amateur's album is art.
\end{objection}

\begin{reply}
Ease bears on the value of the achievement, not on the kind of agency. On Principle~P2 of AES-14-a, artistic value includes achievement relative to the maker's situation (AES-04), and an easy choice is a slight achievement. That selection can be done badly or trivially is no more an objection to it than the existence of trivial collages is an objection to collage. Whether an album is art is a further matter, on AES-01, of presentation and uptake.
\end{reply}

\begin{objection}
The readymade and collage choose whole objects; the photographer chooses among exposures she herself made. The parallel is loose.
\end{objection}

\begin{reply}
The difference strengthens the case. Photographic selection operates on a range the photographer has already shaped by framing, as film editing operates on shots the director has already composed. It is the readymade, where the range is given entirely by the world, that is the limiting case; if selection is agency there, it is agency here \citep{Lopes2016}.
\end{reply}

\section{Conclusion}

Selection is one of the oldest forms of artistic agency in the modern arts. Photography shares it with the readymade, collage and editing, and the objection that would exclude it from photography would exclude it from them. This supports the conclusion of AES-14-a that photography is art through framing and selection, not through accuracy.

\begin{thebibliography}{9}\small
\bibitem[Bordwell(1985)]{Bordwell1985} Bordwell, D. (1985). \emph{Narration in the Fiction Film}. University of Wisconsin Press.
\bibitem[Cartier-Bresson(1952)]{CartierBresson1952} Cartier-Bresson, H. (1952). \emph{The Decisive Moment}. Simon and Schuster.
\bibitem[Danto(1964)]{Danto1964} Danto, A.~C. (1964). The artworld. \emph{The Journal of Philosophy}, 61(19), 571--584.
\bibitem[Lopes(2016)]{Lopes2016} Lopes, D.~M. (2016). \emph{Four Arts of Photography: An Essay in Philosophy}. Wiley-Blackwell.
\end{thebibliography}

\end{document}
